% =============================================================================
%  front.tex —— 扉页、译本说明、体例说明
% =============================================================================
\begin{titlepage}
\centering
\vspace*{2.2cm}
{\LARGE\bfseries 电磁对偶}\\[10pt]
{\LARGE\bfseries 与几何 Langlands 纲领}\\[28pt]
{\large Electric-Magnetic Duality\\[4pt] And The Geometric Langlands Program}\\[34pt]
{\large Anton Kapustin \quad$\cdot$\quad Edward Witten}\\[26pt]
{\normalsize 中文学术译注本}\\[6pt]
{\small 译本附：思路导读、译注、译对报告、插件评估}\\[40pt]
{\small 底本：arXiv:hep-th/0604151v3 \quad 2006 年 4 月（v3：2007 年 1 月），231 页}\\[6pt]
{\small \today}
\vfill
\end{titlepage}

\thispagestyle{empty}
\section*{译本说明}
\addcontentsline{toc}{section}{译本说明}

\noindent\textbf{底本。}本译本的底本是 A.~Kapustin 与 E.~Witten 的
\emph{Electric-Magnetic Duality And The Geometric Langlands Program}
（arXiv:hep-th/0604151v3，231 页，12 节、362 个编号公式（由 348 处公式定义生成）、146 条参考文献）。译文直接从作者提交
arXiv 的 \LaTeX{} 源码（harvmac 格式）逐段生成，而不是从 PDF 文字层转写：公式、编号、
文献编号、图表都来自源码本身，因此可与原书逐式对照。

\medskip
\noindent\textbf{译文。}正文只放译文，不混入解释。译者所加的内容一律放进两类方框：
\begin{itemize}[leftmargin=1.6em,itemsep=2pt,topsep=2pt]
  \item \textbf{本节重点}（蓝框）：每节开头 3--5 条，交代该节要解决的问题、用到的工具、得到的结果；
  \item \textbf{译注}（橙框）：针对具体难点的补白，分「易错点」「含糊点拨」「物理图像」「原文疑误」四类。
\end{itemize}
译注可以整体关闭（\texttt{main.tex} 中 \verb|\showyizhutrue| 改为 \verb|\showyizhufalse|），
关闭后即为纯粹的译文。

\medskip
\noindent\textbf{编号。}编号公式与原文编号一一对应（如原文式 (3.41) 在本译本仍是式 (3.41)），
参考文献沿用原文的 [N] 编号。原文的图形按 \texttt{figures/} 下的图片重新排入，图号与原文一致。

\medskip
\noindent\textbf{原文疑误。}凡译者判断原文有笔误的地方，译文照原文保留，并紧接一条
「译注 [原文疑误]」说明理由与建议改法；不做静默修订。

\medskip
\noindent\textbf{致谢与声明。}原作者的版权、文字与观点归原作者所有；本译本仅供学习与研究使用，
译注部分的文责由译者承担。
